%% Hand-authored circuitikz figure -- NOT generated by maestro2tex, placed by it via
%% the `order` list in configs/anatop_tb_tia.json. Topology transcribed from
%% simarchive/anatop_tb_tia/tia_2026-07-28/Interactive.9/1/<TB>/netlist/input.scs
%% Supplies, the TCBias/LCBias trim buses and En are omitted -- identical in every
%% bench and carrying no measurement.
\begin{figure}[htbp]
  \centering
  \resizebox{0.56\linewidth}{!}{%
\ctikzset{bipoles/length=1.0cm, resistors/scale=0.85, capacitors/scale=0.85, sources/scale=0.85}
\begin{circuitikz}[
    every node/.append style={font=\footnotesize},
    nlab/.style={font=\scriptsize, inner sep=1.5pt},
    opt/.style={font=\scriptsize, text=black!55, align=center, inner sep=2pt},
    railnode/.style={circle, fill=black, inner sep=1.05pt},
    prb/.style={draw, circle, minimum size=9pt, inner sep=0pt},
  ]
  \node[op amp] (A) at (0,0) {};
  \node at (A.center) {\footnotesize A2};
  \draw (A.+) -- (-1.6,-0.3);
  \draw (-1.6,-0.3) to[V, l_=$V_{\mathrm{CM}}$, invert] (-1.6,-2.6) node[ground]{};
  \draw (A.-) -- (-2.8,0.3);
  \draw (-2.8,0.3) -- (-2.8,2.2);
  \draw (-2.8,2.2) -- (-1.9,2.2);
  \node[prb] (P) at (-1.6,2.2) {};
  \node[nlab, above=2pt] at (P.north) {\texttt{IPRB1}};
  \draw (P.east) -- (1.6,2.2) -- (1.6,0);
  \draw (A.out) -- (1.6,0);
  \node[railnode] at (1.6,0){};
  \node[nlab, below right] at (1.7,-0.05){\texttt{V\_OutOL}};
  \draw (1.6,0) -- (2.9,0);
  \draw (2.9,0) to[C, l=$C_L$] (2.9,-2.6) node[ground]{};
  \draw (2.9,0) -- (4.4,0);
  \draw (4.4,0) to[I, l=$i_{\mathrm{load}}$] (4.4,-2.6) node[ground]{};
  \node[opt, anchor=north] at (4.4,-3.0) {\texttt{DriveTB} only};
  \node[opt, anchor=north] at (-1.6,-3.0) {DC; a pulse in\\\texttt{StepUnityFeedbackTB}};
\end{circuitikz}
  }
  \caption[{Unity-feedback testbench for the transimpedance amplifier.}]{Unity-feedback testbench. The amplifier's output is returned to its inverting input through the current probe \texttt{IPRB1} and loaded by \(C_L=\SI{5}{\pico\farad}\); a \texttt{stb} analysis through that probe gives the loop gain and phase, and \texttt{dc} and \texttt{noise} on the same wiring give the offset, tracking range and input-referred noise. Two variants share this schematic: \texttt{DriveTB} adds the swept load current \(i_{\mathrm{load}}\), and \texttt{StepUnityFeedbackTB} replaces the DC source with a pulse. Supplies, the bias trim buses and the enable pin are omitted here and in Figures~\ref{fig:tb-tiafb} and~\ref{fig:tb-slew}.}
  \label{fig:tb-unityfb}
\end{figure}
